这个问题来自"Drinfeld–Sokolov Hierarchies and Diagram Automorphisms of Affine Kac–Moody Algebras".
\begin{conjecture}
Let $V$ be a linear space, $v\in V$, $X,Y,Z\in End(V)$, s.t.

\begin{align*}
    [X,Y]=Z,\,[X,Z]=[Y,Z]=0\\
    X^{p+1}v=0,\,Y^{q+1}v=0
\end{align*}

then $$(X+Y)^{2p+2q+1}v=0$$
\end{conjecture}

Consider the algebra $\mathcal{A}$ generated by $X,Y$ with relations $[X,[X,Y]]=[Y,[X,Y]]=0$, it has a $\mathbb{Z}^2$ grading such that $X$ has degree $(1,0)$ and $Y$ has degree $(0,1)$. Denote degree $(m,n)$ part by $\mathcal{A}_{m,n}$. The degree $(m,n)$ part has basis $Z^iX^{m-i}Y^{n-i}$, $0\le i\le\min(m,n)$.

For convenience we define $Z=[Y,X]$ from now on.

\begin{lemma}
$$Y^nX^m=\sum_{i=0}^{min(m,n)}C_{m,n}^{i}Z^iX^{m-i}Y^{n-i}$$
in which $C_{m,n}^{i}>0$.
\end{lemma}

\begin{proof}
\begin{align*}
    Y^nX^m&=XY^nX^{m-1}+[Y^n,X]X^{m-1}\\
    &=\sum_iC_{m-1,n}^{i}Z^iX^{m-i}Y^{n-i}+nZY^{n-1}X^{m-1}\\
    &=\sum_iC_{m-1,n}^{i}Z^iX^{m-i}Y^{n-i}+\sum_i nC_{m-1,n-1}^{i-1}Z^iX^{m-i}Y^{n-i}\\
    &=\sum_i(C_{m-1,n}^{i}+nC_{m-1,n-1}^{i-1})Z^iX^{m-i}Y^{n-i}\\
\end{align*}
So $C_{m,n}^{i}$ are sum of positive numbers, thus always positive.
\end{proof}

\begin{remark}
    $$C_{m,n}^{i}=i!\binom{m}{i}\binom{n}{i}$$
    
    Consider the morphism from $\mathcal{A}$ to Weyl algebra by $X \mapsto x$, $Y \mapsto \partial$. The morphism is injective when restricted on $\mathcal{A}_{m,n}$. Then one can apply higher Leibniz rule.
\end{remark}
\begin{remark}
We only need $C_{m,n}^{n}\not =0$ when $m\ge n$.
\end{remark}

\begin{proof}[Proof of conjecture 1]
It is suffice to prove that if $m+n=2p+2q+1$ and $M\in \mathcal{A}_{m,n}$ then $Mv=0$. W.L.O.G assume $n\ge p+q+1$.

Let $M_i=Z^iX^{m-i}Y^{n-i}$ for $0\le i\le m$ and $M^{\prime}_i=Z^i X^{m-i-p-1} Y^{n-i} X^{p+1}$ for $0\le i\le m-p-1$. So $M^{\prime}_i v=0$, and $M_i v=0$ if $i\le n-q-1$. 

Denote span of $\{M^{\prime}_i\}\bigcup\{M_i| i\le n-q-1\}$ by $W$ we need to show $W=\mathcal{A}_{m,n}$, namely $M_i\in W$ for all $i$. Expand $M^{\prime}_i$ by $M_i$, we get: 

$$M^{\prime}_i=\sum_{j=0}^{p+1}C_{n-i,p+1}^{j}M_{i+j}$$

Note that $m\le p+q$, then $i\le m-p-1 \le q-1$, so $n-i\ge p+1$, thus $C_{n-i,p+1}^{j}\not =0$. 

So we can perform induction: assume $M_k\in W$ for all $k < j$, in which $j>n-q$. Expand $M^{\prime}_{j-p-1}$, we get linear combination of $M_{j-p-1},\dots,M_{j}$. So $M_{j}\in S$. And we can reach $j=m$ by expanding $M_{m-p-1}$.


%\color{red}{$M^{\prime}_i$ are in linear span of $M_j$ with $i\le j\le(i+p+1)$. The coefficients should be a upper triangular matrix with non-vanishing diagonal. As $M^{\prime}_i v=0$, we have $M_iv=0$.}
\end{proof}
